\label{sec:conclusion}

We set out to build a better memory for autonomous agents on associative-memory foundations.
That thesis is refuted here by our own measurements, at adequate power, with the mechanism
identified: every departure from the dense-retrieval corner that changes only how an answer
is read off the energy landscape fails on agent text, because none of them changes the
landscape. The one departure that survives changes the \emph{cue} instead, and it is query
expansion rather than attractor dynamics.

What replaced it is narrower and better supported. Dense retrieval is exactly the
$(T{=}1,\beta\!\to\!\infty,b{=}0,R{=}I)$ corner of the modern Hopfield update --- exactly
so at $k{=}1$, and at $k>1$ with the shortlist read off the cue --- which makes
the standard baseline a configuration of the method under test rather than a rival outside
it, and converts an ill-posed question into a finite one. Within that frame, one side of
the sign is predictable before the experiment: associative descent hurts monotonically when
the cue is not a degraded instance of a stored pattern, and that half has held everywhere
we have looked. The other half has no surviving instance. We nominated three tests against
ourselves and all three came back negative: satisfying the premise on masked text gives no
winning $\beta$ in $45$ cells, two image targets instead of one costs $-0.245$, and the
single-target image result we had treated as the confirming case turns out to be a ten-way
classification measured on cues whose source was never stored, which an explicit $k$-NN
class vote performs better than descent does ($-0.036$, \ci{-0.043}{-0.028}). On the
premise's own best case --- occluded copies of stored images, scored by identity --- descent
is $-0.003$ \ci{-0.008}{+0.000}. The premise is therefore necessary and not sufficient, and
we withdraw the image win rather than keep it under a weaker description
(\S\ref{sec:occlusion}). Of eight mechanisms beyond the corner, one survives---a
nonparametric predictor over the agent's own observed transitions, which retrieves by
succession rather than resemblance. It beats cosine on one-hop association in four
measurements across three corpora, including $+0.043$ Recall@1 \ci{+0.036}{+0.050} on
\textsc{LoCoMo} against a pooled bank of $5{,}882$ turns and $+0.016$ \ci{+0.014}{+0.018}
on \textsc{MSC} against $30{,}320$. The same experiments bound it: at two hops the effect is
$-0.002$ on the first corpus and $+0.005$ on the second, one eighteenth of the $+0.185$
two-hop result of our own that we withdraw. Association helps one step ahead and essentially
stops there.

\paragraph{One thing beats the corner, and it is not the thing we were looking for.} We
could not tell from \textsc{LoCoMo} alone whether query expansion fails as a mechanism or
fails on that corpus, so we ran six more. It fails on the corpus. Fold a few retrieved
memories back into the cue and retrieve again, and on multi-hop questions whose second
document is reachable only \emph{from the first} the gain is $+0.023$ \ci{+0.010}{+0.037}
and $+0.171$ \ci{+0.118}{+0.219}; on questions in the same corpora whose documents are both
reachable from the question, selection declines to expand in $10$ of $10$ folds and the
method correctly does nothing. Random modes lose to real ones by $+0.13$ to $+0.39$, so what
pays is which memories enter the cue. Our own pre-registered explanation for when this
should happen---more targets---comes out reversed on four corpora, and the predictor that
survives is the shape of the hop. We hold the claim at arm's length nonetheless: it needs
per-corpus tuning, it needs a question shape we cannot yet detect without labels, and on the
two conversational corpora that motivate this paper it is null, where a lexical hybrid gains
$+0.156$. It is the one place where the thesis we began with was too pessimistic, and it is
still not a better agent memory (\S\ref{sec:corpora}).

We also report what the protocol was worth, in the only currency that makes the argument:
seven errors, each a reasonable default, each worth a conclusion. Two produced false
negatives, which were materially harder to detect than the false positives. Three of our own
recorded claims were overturned by reimplementation rather than by review---and the one that
had been our headline was settled only by the public-benchmark run that external validity
demanded anyway. The seventh is the one a reader should weigh most heavily, because it
survived every check we had: an \emph{unmatched baseline}. The image comparison was between
a 1-nearest-neighbour classifier and a soft-$k$-nearest-neighbour one, on a class metric,
and no amount of split hygiene, pre-registration or bootstrap discipline detects that. Only
naming the operation the method actually performs does.

The constructive proposal we take from this is structural rather than empirical. Build the
baseline as a corner of the method, give every added component a null setting that reproduces
the configuration beneath it, and assert those equivalences in continuous integration. Four
of the seven errors above surfaced as violations of a known equivalence rather than as
anomalous numbers. In a literature where the operator under test and the operator it is
compared against coincide at the default operating point, that is not a convenience. It is
the difference between measuring a method and measuring a bug.
